\documentclass[12pt,a4paper]{article}

\usepackage[a4paper,margin=1in]{geometry}
\usepackage{graphicx}
\usepackage{booktabs}
\usepackage{amsmath,amssymb}
\usepackage{float}
\usepackage{hyperref}
\usepackage{xcolor}
\usepackage{listings}
\usepackage{enumitem}
\usepackage{fancyhdr}
\usepackage{longtable}

\hypersetup{colorlinks=true,linkcolor=blue,urlcolor=blue}

\pagestyle{fancy}
\fancyhf{}
\lhead{ICS1512}
\rhead{Machine Learning Algorithms Laboratory}
\cfoot{\thepage}

\lstset{
language=Python,
basicstyle=\ttfamily\small,
numbers=left,
frame=single,
breaklines=true,
showstringspaces=false
}

\begin{document}

\begin{tabular}{|p{4cm}|p{11cm}|}
\hline
Experiment No. & 3 \\ \hline
Experiment Title & Regression Analysis using Linear and Regularized Models
\footnote{\textbf{GitHub Repository: }\url{https://github.com/Barath-j593/ML-assignments}}\\ \hline
Student Name & NavinKumar S \\ \hline
Register Number & 3122247001039\\ \hline
Faculty & Dr. Poreddy Ajay Kumar Reddy \\ \hline
Submission Date & 02/08/26 \\ \hline
\end{tabular}

\section{Objective}
To implement Linear, Ridge, Lasso, and Elastic Net regression models to predict a continuous
target variable -- the requested loan amount -- tune regularization hyperparameters via
GridSearchCV, evaluate performance using multiple regression metrics, visualize model
behaviour, and analyze overfitting, underfitting, and bias--variance characteristics.

\section{Problem Statement}
Given the Kaggle \textit{Predict Loan Amount Data} dataset of 20,000 customer loan
applications with demographic, financial, employment, and property-related features, the
problem is regression: predict the continuous \texttt{Loan Amount Request (USD)} for each
applicant. The input is a mixed categorical/numerical feature vector per applicant; the
output is a continuous dollar value. The objective is to determine which regularization
strategy (if any) generalizes best, and to characterize the bias--variance behaviour of each
model on this dataset.

\section{Dataset Description}
\begin{longtable}{|p{6cm}|p{8cm}|}
\hline
Dataset Name & Predict Loan Amount Data \\ \hline
Dataset Source & Kaggle -- \textit{Predict Loan Amount Data} \\ \hline
Number of Samples & 20,000 \\ \hline
Number of Features & 19 used (10 categorical, 9 numerical) -- expands to 45 after one-hot encoding \\ \hline
Number of Classes & N/A -- regression target (continuous), range 6{,}185--576{,}336 USD \\ \hline
Missing Values & Gender (31), Income (750), Income Stability (813), Type of Employment (4{,}689), Current Loan Expenses (83), Dependents (1{,}142), Credit Score (743), Has Active Credit Card (1{,}076), Property Age (892), Property Location (160), plus disguised ``?'' placeholders in Property Price (168) and Co-Applicant (77) \\ \hline
Train-Test Split & 80:20 (16,000 train / 4,000 test), \texttt{random\_state=42} \\ \hline
\end{longtable}

\section{Exploratory Data Analysis}
Include all relevant visualizations.
\begin{itemize}
\item Dataset overview
\item Statistical summary
\item Missing value analysis
\item Class distribution (not applicable -- regression target)
\item Correlation matrix (feature-vs-target correlations reported below)
\item Feature distribution
\end{itemize}

\subsection*{Target Distribution}
\begin{figure}[H]
\centering
\includegraphics[width=0.6\linewidth]{target_distribution.png}
\caption{Target variable (Loan Amount Request) distribution}
\end{figure}
\textbf{Inference:} The requested loan amount is right-skewed (skewness $\approx 1.28$),
with most requests concentrated below 150{,}000 USD and a tail extending to over 500{,}000
USD. With 20,000 samples, the extreme tail has limited influence on the overall shape
compared to a smaller dataset.

\subsection*{Feature Distribution / Correlation with Target}
\begin{figure}[H]
\centering
\includegraphics[width=0.55\linewidth]{scatter_Property.png}
\caption{Property Price vs. Loan Amount Request}
\end{figure}
\textbf{Inference:} Property Price shows an extremely strong, almost linear positive
relationship with the requested loan amount ($r \approx 0.961$) -- by far the strongest
relationship of any feature, and the dominant driver of every model's predictions.

\begin{figure}[H]
\centering
\includegraphics[width=0.55\linewidth]{scatter_Income.png}
\caption{Income (USD) vs. Loan Amount Request}
\end{figure}
\textbf{Inference:} Income shows a much weaker positive relationship with loan amount
($r \approx 0.364$) than Property Price -- applicants with similar incomes still request
widely varying loan amounts, indicating income alone is a poor predictor without the
property-value context.

\section{Data Preprocessing}
Explain:
\begin{itemize}
\item \textbf{Missing value handling:} \texttt{Property Price} and \texttt{Co-Applicant}
encode missing data as the literal string ``?'' rather than a blank cell; both were converted
to \texttt{NaN} and then to numeric type. All categorical columns with missing values were
then imputed with their mode, and all numerical columns with missing values were imputed with
their median. \texttt{Customer ID}, \texttt{Name}, and \texttt{Property ID} were dropped as
unique identifiers with no generalizable signal.
\item \textbf{Encoding:} All ten categorical features (Gender, Income Stability, Profession,
Type of Employment, Location, Expense Type 1, Expense Type 2, Has Active Credit Card,
Property Type, Property Location) were one-hot encoded (\texttt{drop="first"}), expanding 19
raw input features into 45 encoded features.
\item \textbf{Feature scaling:} Age, Income, Current Loan Expenses, Dependents, Credit Score,
No. of Defaults, Property Age, Co-Applicant, and Property Price were standardized (zero mean,
unit variance) using \texttt{StandardScaler}, fit only on the training split.
\item \textbf{Train-test split:} 80:20 (16,000 train / 4,000 test rows), implemented inside a
single \texttt{Pipeline} (preprocessing $\rightarrow$ model) to prevent data leakage across
cross-validation folds and the held-out test set.
\item \textbf{Feature engineering:} Not applied in this experiment -- the 19 raw features
were used directly after encoding/scaling.
\end{itemize}

\section{Machine Learning Algorithm}
Explain:
\begin{itemize}
\item \textbf{Working principle:} Linear Regression fits a linear combination of input
features to minimize squared error against the continuous target. Ridge Regression adds an
$L_2$ penalty on coefficient magnitudes to shrink them toward zero without eliminating any.
Lasso Regression adds an $L_1$ penalty that can shrink coefficients exactly to zero,
performing implicit feature selection. Elastic Net linearly combines both $L_1$ and $L_2$
penalties, controlled by a mixing parameter (l1\_ratio).
\item \textbf{Mathematical formulation:}
$$\hat{y} = w^\top x + b, \qquad
J(w) = \underbrace{\|y - Xw\|_2^2}_{\text{OLS loss}} + \underbrace{\alpha\|w\|_2^2}_{\text{Ridge}}
\ \text{or}\ \underbrace{\alpha\|w\|_1}_{\text{Lasso}}
\ \text{or}\ \underbrace{\alpha\left(r\|w\|_1 + \tfrac{1-r}{2}\|w\|_2^2\right)}_{\text{Elastic Net}}$$
where $\alpha$ controls overall regularization strength and $r$ (l1\_ratio) controls the
$L_1$/$L_2$ mix in Elastic Net.
\item \textbf{Advantages:} Linear Regression is simple, fast, and fully interpretable.
Ridge handles multicollinearity gracefully by shrinking correlated coefficients together.
Lasso performs automatic feature selection, useful when many features are irrelevant.
Elastic Net balances both behaviours when the number/correlation structure of useful features
is uncertain.
\item \textbf{Limitations:} All four are linear models and cannot capture non-linear
feature--target relationships. Ridge never produces a sparse (zero-coefficient) model even
when many features are irrelevant. Lasso can behave erratically with strongly correlated
features (arbitrarily picking one over another). Elastic Net introduces a second
hyperparameter to tune, increasing search cost.
\end{itemize}

\section{Experimental Setup}
\begin{tabular}{ll}
\toprule
Python Version & \\
Scikit-Learn Version & \\
IDE & Jupyter Notebook \\
Operating System & \\
Hardware & \\
\bottomrule
\end{tabular}

\section{Hyperparameter Selection}

\begin{tabular}{ll}
\toprule
Parameter & Value\\
\midrule
Ridge -- best $\alpha$ (GridSearchCV) & 10 \\
Ridge -- best CV $R^2$ & 0.9302 \\
Lasso -- best $\alpha$ (GridSearchCV) & 10 \\
Lasso -- best CV $R^2$ & 0.9303 \\
Elastic Net -- best $\alpha$ (GridSearchCV) & 0.01 \\
Elastic Net -- best l1\_ratio (GridSearchCV) & 0.8 \\
Elastic Net -- best CV $R^2$ & 0.9302 \\
\bottomrule
\end{tabular}

\textbf{Note:} Search spaces followed the assignment specification: Ridge
$\alpha \in \{0.01, 0.1, 1, 10, 100\}$; Lasso $\alpha \in \{0.001, 0.01, 0.1, 1, 10\}$;
Elastic Net $\alpha \in \{0.01, 0.1, 1, 10\}$, l1\_ratio $\in \{0.2, 0.5, 0.8\}$, all searched
via 5-fold \texttt{GridSearchCV} optimizing $R^2$.

\section{Results}

\begin{tabular}{lc}
\toprule
Metric (Regression) & Value (Best Model: Lasso Regression, Test Set)\\
\midrule
MAE & 10{,}890.31 \\
MSE & 232{,}800{,}718 \\
RMSE & 15{,}257.81 \\
R\textsuperscript{2} & 0.9356 \\
Training Time (s) & 0.0928 \\
\bottomrule
\end{tabular}

\textbf{Note:} The template's classification metrics (Accuracy/Precision/Recall/ROC-AUC) do
not apply to this regression task and have been replaced with the standard regression
metrics specified in the assignment (MAE, MSE, RMSE, $R^2$, Training Time). All four models
performed within 0.01\% $R^2$ of each other on the test set (Linear 0.9355, Ridge 0.9355,
Lasso 0.9356, Elastic Net 0.9356).

\section{Result Visualization}

Include all graphs generated during the experiment.

\begin{figure}[H]
\centering
\includegraphics[width=0.55\linewidth]{predicted_vs_actual.png}
\caption{Predicted vs. actual Loan Amount Request -- best model (Lasso Regression)}
\end{figure}
\textbf{Inference:} Points cluster tightly along the ideal $y=x$ line across almost the
entire range, only fanning out slightly at the highest loan amounts -- visually confirming
the high test $R^2$ ($\approx 0.936$) and showing the model generalizes well.

\begin{figure}[H]
\centering
\includegraphics[width=0.55\linewidth]{residual_plot.png}
\caption{Residual plot -- best model (Lasso Regression)}
\end{figure}
\textbf{Inference:} Residuals are centered close to zero (mean $\approx -309$, small
relative to the target's scale) and show only mild fanning at higher predicted values --
a much weaker heteroscedasticity pattern than the target's skew alone would suggest, since
Property Price's near-linear relationship with the target absorbs most of that structure.

\begin{figure}[H]
\centering
\includegraphics[width=0.55\linewidth]{train_val_error.png}
\caption{Learning curve (training vs. validation RMSE) -- best model (Lasso Regression)}
\end{figure}
\textbf{Inference:} Training RMSE ($\approx$15{,}700--16{,}500) and validation RMSE
($\approx$15{,}800--16{,}200) stay close together at every training size and both decrease as
more data is added, converging toward a shared, low-error plateau -- the signature of a
well-fit model with neither severe overfitting nor severe underfitting.

\begin{figure}[H]
\centering
\includegraphics[width=0.85\linewidth]{coefficient_comparison.png}
\caption{Feature importance -- coefficient comparison across models (top 8 features by |Linear coefficient|)}
\end{figure}
\textbf{Inference:} Property Price dominates every model's coefficients by an order of
magnitude over any other feature (Linear: 51{,}583 vs. the next largest at 7{,}727 for
Current Loan Expenses). Lasso zeroes out three of the remaining top-8 features entirely
(Profession\_Unemployed, Type of Employment\_Realty agents, Type of Employment\_IT staff),
while Ridge and Elastic Net shrink but retain all of them -- the expected Lasso-vs-Ridge
sparsity contrast.

\section{Discussion}
Unlike a weaker-signal regression problem, this dataset -- with Property Price as a feature
-- produces a strong, well-behaved fit: all four models achieve test $R^2 \approx 0.935$--
$0.936$. Because the unregularized Linear Regression baseline was already close to the
regularized models' performance, regularization strength was not the binding constraint here
-- there was very little overfitting for $L_1$/$L_2$ penalties to correct. Lasso achieved the
best (marginally) test performance by trimming a handful of low-signal categorical dummy
variables without disturbing the dominant numeric predictors. A limitation of this experiment
is that with one feature (Property Price) so overwhelmingly dominant, the experiment offers
limited insight into how these models would behave on a dataset with more diffuse,
comparably-weighted feature signal; a useful follow-up would be to refit after removing
Property Price to observe whether regularization becomes more consequential.

\section{Conclusion}
For the Predict Loan Amount Data dataset, Lasso Regression (tuned to $\alpha=10$) is
marginally the best-performing model on both cross-validation ($R^2=0.9303$) and the held-out
test set (RMSE $=15{,}257.81$, $R^2=0.9356$), though all four models are within 0.01\%
$R^2$ of one another. Property Price alone accounts for the overwhelming majority of
predictive power. The learning curve shows converging, low training/validation error with no
meaningful gap, indicating this experiment sits in a low-bias, low-variance regime where
additional regularization tuning offers negligible practical benefit over plain Linear
Regression.

\section{References}
\begin{enumerate}[label={[\arabic*]}]
\item Scikit-learn Developers, ``Linear Models,'' \url{https://scikit-learn.org/stable/modules/linear_model.html}
\item Scikit-learn Developers, ``Tuning the hyper-parameters of an estimator,'' \url{https://scikit-learn.org/stable/modules/grid_search.html}
\item Kaggle, ``Predict Loan Amount Data,'' \url{https://www.kaggle.com/datasets}
\end{enumerate}

\end{document}